\section{Contoh Program}
\label{sec:contoh-bab9}

Bagian ini menyajikan dua program utuh. Program pertama memakai
seluruh direktif data, sedangkan program kedua memakai \dir{EQU}
sebagai konstanta sekaligus memperlihatkan apa yang terjadi bila
direktif itu didefinisikan ulang.

\subsection{Contoh 1: Mendemonstrasikan Semua Direktif Data}

Program pada Kode~\ref{lst:direktif-data} memakai \dir{DB},
\dir{DW}, \dir{DD}, \dir{DUP}, dan penanda \code{?}, lalu mencetak
sebagian isi memori tersebut ke layar.

\lstinputlisting[
    language=MASM,
    caption={Program \file{direktif-data.asm}: direktif data dan
    isi memorinya.},
    label={lst:direktif-data}
]{\CodeDir/direktif-data.asm}

Bagian \dir{.DATA} pada program tersebut dapat dibaca sebagai
berikut.

\begin{itemize}
    \item \lstasm|sapaan DB 'Halo, assembly!', 0Dh, 0Ah, '$'|
          menyediakan \dec{18} byte: \dec{15} byte isi string,
          satu byte \code{0Dh} dan satu byte \code{0Ah} untuk
          pindah baris, serta \code{'\$'} sebagai penanda akhir
          untuk layanan DOS fungsi \hex{09}.
    \item \lstasm|panjang EQU $ - sapaan| meminta assembler
          menghitung panjang itu sendiri. Nilainya terlihat pada
          berkas \file{.lst} (Kode~\ref{lst:listing-lst}).
    \item \lstasm|kata DW 4241h, 4443h| menyediakan empat byte.
          Karena urutan \eng{little-endian}, byte yang tersimpan
          berurutan \hex{41}, \hex{42}, \hex{43}, \hex{44}.
    \item \lstasm|nilai DD 12345678h| menyediakan empat byte tanpa
          dipakai program, tetapi tetap menggeser \eng{offset}
          \code{bintang} dan \code{buffer}.
    \item \lstasm|bintang DB 4 DUP('*')| menunjukkan \dir{DUP} yang
          berisi nilai awal, sedangkan \lstasm|buffer DB 8 DUP(?)|
          hanya menyediakan ruang kosong.
\end{itemize}

Bagian \dir{.CODE} mencetak sapaan, lalu membaca \code{kata}
\emph{byte per byte} dengan \lstasm|BYTE PTR [kata+1]|. Cara
pengalamatan bertanda kurung siku dibahas pada
Bab~\ref{chap:mode-pengalamatan}.

\begin{session}[caption={Keluaran program \file{direktif-data.asm}.}, label={lst:keluaran-direktif-data}]
Halo, assembly!
ABCD****
\end{session}

Baris kedua keluaran tersebut membuktikan urutan
\eng{little-endian}: data \hex{4241} dan \hex{4443} terbaca di
memori sebagai \code{ABCD}.

\subsection{Contoh 2: Konstanta \dir{EQU}}

Program pada Kode~\ref{lst:equ-konstanta} mengganti setiap nomor
fungsi DOS dengan nama yang bermakna. Nilainya tidak berubah
dibandingkan penulisan langsung; yang berubah hanyalah keterbacaan
program.

\lstinputlisting[
    language=MASM,
    caption={Program \file{equ-konstanta.asm}: konstanta simbolik
    dengan \dir{EQU}.},
    label={lst:equ-konstanta}
]{\CodeDir/equ-konstanta.asm}

\begin{session}[caption={Keluaran program \file{equ-konstanta.asm}.}, label={lst:keluaran-equ-konstanta}]
Konstanta EQU
X = 7
\end{session}

\subsection{Definisi Ulang \dir{EQU} Dibandingkan dengan \code{=}}

Kode~\ref{lst:equ-ulang-gagal} menimbulkan kesalahan perakitan yang
sama seperti pada Kode~\ref{lst:equ-salah}.

\begin{lstlisting}[
    language=MASM,
    caption={Baris yang menimbulkan kesalahan perakitan.},
    label={lst:equ-ulang-gagal}
]
NILAI   EQU 5           ; definisi pertama
NILAI   EQU 9           ; gagal: Symbol already defined
\end{lstlisting}

Bila kata kunci \dir{EQU} diganti dengan \code{=}, kedua baris itu
menjadi sah. Program berikut mencetak dua angka berbeda dari satu
nama simbol.

\begin{lstlisting}[
    language=MASM,
    caption={\code{=} boleh didefinisikan ulang; program mencetak
    \code{59}.},
    label={lst:sama-dengan-ulang}
]
.MODEL SMALL
.STACK 100h

.DATA
    awalan  DB 'Nilai = ', '$'

.CODE
START:
    MOV AX, @DATA
    MOV DS, AX

    MOV AH, 09h             ; cetak awalan
    LEA DX, awalan
    INT 21h

NILAI   = 5                 ; definisi pertama
    MOV AH, 02h             ; cetak hasil '0' + 5 = '5'
    MOV DL, '0' + NILAI
    INT 21h

NILAI   = 9                 ; definisi ulang: sah
    MOV AH, 02h             ; cetak hasil '0' + 9 = '9'
    MOV DL, '0' + NILAI
    INT 21h

    MOV AH, 4Ch
    MOV AL, 00h
    INT 21h

END START
\end{lstlisting}

\begin{session}[caption={Keluaran program pada Kode~\ref{lst:sama-dengan-ulang}.}]
Nilai = 59
\end{session}

Kesimpulannya sederhana: pakai \dir{EQU} untuk nilai yang
\emph{tidak akan pernah berubah}, seperti nomor fungsi DOS dan
ukuran tetap, dan pakai \code{=} untuk nilai yang memang harus
diganti di tengah berkas.
